\paragraph{Pregunta para investigar}

\begin{enumerate}[label=P\arabic*]
	\item\label{item:piso}
		?`Qu\'e es la \emph{funci\'on piso} $\piso a$?
		Calcular $\piso{\sqrt 5}$, $\piso{\sqrt{-2}}$,
		$\piso{\frac{\sqrt 7} 2}$, $\piso{2+\sqrt 3}$.
	\seti
\end{enumerate}
%

\paragraph{Una pregunta para explorar}
El sistema $\polinomios{\Enterosmod[3]}$ se parece mucho a $\Enteros$.
\begin{enumerate}[label=P\arabic*]
	\conti
	\item\label{item:cuerpos}
		Consider\'a el polinomio $\indet^2+1$ en
		$\polinomios{\Enterosmod[3]}$ y describ\'{\i} todos
		los elementos del sistema matem\'atico
		$F=\big(\polinomios{\Enterosmod[3]}\big)_{\indet^2+1}$.
	\seti
\end{enumerate}
%
En el problema \ref{item:cuerpos}, el sistema $F$ es el sistema de
restos de la divisi\'on de polinomios en $\polinomios{\Enterosmod[3]}$
por el polinomio $\indet^2+1$.
\begin{enumerate}[label=P\arabic*]
	\conti
	\item\label{item:cuerpos:producto}
		Calcul\'a el producto $(\indet+1)\,(\indet+2)$ en $F$.
	\item\label{item:cuerpos:inversos}
		Calcul\'a los siguientes inversos en $F$:
		\begin{displaymath}
			\frac 1 \indet\quad\text{,}\quad
			\frac 1{\indet^2+2}\quad\text{,}\quad
			\frac 1{\indet^2+2\indet+2}\text{ .}
		\end{displaymath}
		%
	\seti
\end{enumerate}
%

\paragraph{Propiedades de los n\'umeros enteros}
Demostr\'a o refut\'a cada una de las siguientes propiedades.
En caso de que una afirmaci\'on sea falsa en
general, enmendarla, si es posible.
\begin{enumerate}[label=P\arabic*]
	\conti
	\item\label{item:cancelatividad}
		Si $ac=ab$, entonces $a=b$ ?`Es cierta en $\Enteros$?
		?`Es cierta en $\polinomios[\raizcuarta] \Enteros$?
		?`Lo es en $\Enterosmod[m]$?
	\item\label{item:naturales:clausura}
		Si $a,b\in\Enteros$, entonces el producto $ab\in\Naturales$,
		si y s\'olo si $a$ y $b$ son \emph{ambos} naturales.
	\item\label{item:cuadrados}
		Si $a\neq 0$, entonces $a^2>0$ ?`Es esta afirmaci\'on
		cierta en $\Enteros$? ?`Es cierta en otros sistemas
		que conozcas?
	\item\label{item:division}
		Si $a,b\in\Enteros$, entonces existen enteros
		$q$ y $r$ tales que $a=bq+r$ y $0\leq r<|b|$.
	\item\label{item:reducibilidad}
		Todo entero $>1$ posee alg\'un divisor primo.
	\item\label{item:discreto}
		No existen enteros entre $0$ y $1$.
	\item\label{item:mersenne}
		Si $a$ y $n$ son enteros positivos, $>1$ y
		$a^n-1$ es primo, entonces necesariamente
		$a=2$ y $n$ debe ser primo.
	\seti
\end{enumerate}
%

\paragraph{Problemas num\'ericos}
El algoritmo de Euclides puede tener aplicaciones inesperadas.
\begin{enumerate}[label=P\arabic*]
	\conti
	\item\label{item:mcd}
		Usando el algoritmo de Euclides, calcul\'a
		el m\'aximo com\'un divisor
		$\mcd{7469,2464}$.
	\item\label{item:diofantica}
		Encontr\'a una soluci\'on a la ecuaci\'on
		$2464x+7469y=77$ con $x$ e $y$ enteros.
	\item\label{item:inversos}
		Hall\'a las soluciones a la ecuaci\'on
		$2464x=11$ en $\Enterosmod[7469]$.
		Hall\'a las soluciones a $2464x=9$ en $\Enterosmod[7469]$.
	\item\label{item:logaritmos}
		Hac\'e una tabla de logaritmos en $\Unidadesmod[29]$
		y us\'a tu tabla para encontrar todas las soluciones
		a las siguientes ecuaciones en $\Enterosmod[29]$:
		\begin{displaymath}
			13x^3\,=\,21\quad\text{;}\quad
			x^4\,=\,7\quad\text{;}\quad
			x^7\,=\,18\quad\text{;}\quad
			x^7\,=\,1\text{ .}
		\end{displaymath}
		%
	\item\label{item:primitivas}
		Hall\'a todas las ra\'{\i}ces primitivas en
		$\Unidadesmod[19]$.
	\seti
\end{enumerate}
%

\paragraph{La t\'ecnica de generalizaci\'on}
Veamos si podemos generalizar el algoritmo de Euclides a otras situaciones.
\begin{enumerate}[label=P\arabic*]
	\conti
	\item\label{item:generalizacion}
		Describ\'{\i} el algoritmo de Euclides de tal manera
		que sea posible aplicar el proceso descrito en los
		problemas P6, P7 y P8 de la Lista~\#2 a $\sqrt 3$ y
		a otros irracionales.
	\seti
\end{enumerate}
%

\paragraph{El arte de contar}
Escribimos $\indice A$ para indicar la cantidad de elementos de un
conjunto.
\begin{enumerate}[label=P\arabic*]
	\conti
	\item\label{item:conjuntos}
		Demostr\'a las siguientes afirmaciones ($A$, $B$,
		$C$ y $D$ son conjuntos):
		\begin{enumerate}[label=(\alph*)]
			\item
				\begin{math}
					\indice{A\cup B}=
					\indice A+\indice B
						-\indice{A\cap B}
				\end{math};
			\item
				\begin{math}
					\indice{A\cup B\cup C}=
					\indice A+\indice B+\indice C
						-\indice{A\cap B}
						-\indice{B\cap C}
						-\indice{C\cap A}
					+\indice{A\cap B\cap C}
				\end{math};
			\item
				\begin{math}
					\indice{A\cup B\cup C\cup D}=
					\text{?`?}
				\end{math}
		\end{enumerate}
		%
	\seti
\end{enumerate}
%

\paragraph{Ingenio}
\begin{enumerate}[label=P\arabic*]
	\conti
	\item\label{item:factorial}
		?`Cu\'antos ceros aparecen al final de la expansi\'on
		decimal de $\factorial{100}$ (cien factorial)?
		?`Cu\'antos ceros aparecen al final de la expansi\'on
		\emph{binaria} de $\factorial{100}$? En ambos casos,
		la pregunta hace referencia a $\enbase{10}{100}$.
\end{enumerate}
%

